% Options for packages loaded elsewhere
\PassOptionsToPackage{unicode}{hyperref}
\PassOptionsToPackage{hyphens}{url}
\PassOptionsToPackage{dvipsnames,svgnames,x11names}{xcolor}
\documentclass[
  11pt,
  letterpaper,
]{article}
\usepackage{xcolor}
\usepackage[margin=1in]{geometry}
\usepackage{amsmath,amssymb}
\setcounter{secnumdepth}{-\maxdimen} % remove section numbering
% Classic Computer Modern; no system-font or Unicode-math fallback.
\usepackage{iftex}
\usepackage[OT1]{fontenc}
\renewcommand{\rmdefault}{cmr}
\renewcommand{\sfdefault}{cmss}
\renewcommand{\ttdefault}{cmtt}

\ifPDFTeX\else
  % xetex/luatex font selection
\fi
% Use upquote if available, for straight quotes in verbatim environments
\IfFileExists{upquote.sty}{\usepackage{upquote}}{}
\IfFileExists{microtype.sty}{% use microtype if available
  \usepackage[]{microtype}
  \UseMicrotypeSet[protrusion]{basicmath} % disable protrusion for tt fonts
}{}
\makeatletter
\@ifundefined{KOMAClassName}{% if non-KOMA class
  \IfFileExists{parskip.sty}{%
    \usepackage{parskip}
  }{% else
    \setlength{\parindent}{0pt}
    \setlength{\parskip}{6pt plus 2pt minus 1pt}}
}{% if KOMA class
  \KOMAoptions{parskip=half}}
\makeatother
\setlength{\emergencystretch}{3em} % prevent overfull lines
\providecommand{\tightlist}{%
  \setlength{\itemsep}{0pt}\setlength{\parskip}{0pt}}
\usepackage{microtype}
\usepackage{amssymb}
\setlength{\emergencystretch}{3em}
\allowdisplaybreaks[1]
\usepackage{bookmark}
\IfFileExists{xurl.sty}{\usepackage{xurl}}{} % add URL line breaks if available
\urlstyle{same}
% fallback for those not using the hyperref driver hyperxmp:
\makeatletter
\@ifundefined{xmpquote}{\newcommand{\xmpquote}[1]{#1}}{}
\makeatother
\hypersetup{
  pdftitle={Further consequences of particular-code knowledge},
  colorlinks=true,
  linkcolor={Maroon},
  filecolor={Maroon},
  citecolor={Blue},
  urlcolor={blue},
  pdfcreator={LaTeX via pandoc}}

\title{Further consequences of particular-code knowledge}
\usepackage{etoolbox}
\makeatletter
\providecommand{\subtitle}[1]{% add subtitle to \maketitle
  \apptocmd{\@title}{\par {\large #1 \par}}{}{}
}
\makeatother
\subtitle{Supplement to the successor-chain paper}
\author{}
\date{20 October 2026}

\begin{document}
\maketitle
\begin{abstract}
We record consequences of the argument for Alexander's Open Problem
3.2.1 that are useful separately from the final impossibility theorem.
Full particular-code knowledge propagates effective enumerability from
one machine to all machines at finite successor distance. The
impossibility theorem excludes computable infinite endorsement paths,
and therefore finite endorsement cycles, under its arithmetic and
truthfulness hypotheses. A single endorsed arithmetic truth column
yields an explicit arithmetic sentence separating two adjacent theories.
Finally, a true arithmetical condition of complexity at most
\(\Sigma_3\) cannot support a conditional code in the presence of a
specified inclusion of arithmetic consequences. All statements below
have proofs. References to the companion paper are to files within this
repository; no other research notes are needed.
\end{abstract}

\section{1. Setting and relation to the main
paper}\label{setting-and-relation-to-the-main-paper}

The companion paper is
\href{../successor-knowers-3-2-1.md}{\emph{Successor truthfulness and
particular-code knowledge}}, referred to below as \textbf{{[}Main{]}}.
It defines Alexander's standard canonical semantics, the code formula,
the partial truth predicate and the witnessed arithmetic theories. Here
we repeat the conventions needed to state the additional results.

Work in a recursively presented language with arithmetic and modal
operators \(K_i\). In a common standard canonical interpretation \(M\),
\[
M\models K_i\varphi[s]
\quad\Longleftrightarrow\quad
T_i\vdash\varphi^s.
\tag{S1}
\] All consequences of each theory under consideration are assumed true
in \(M\). Define \[
\operatorname{Code}_{i,\varphi}(e)
:=\forall x\bigl(K_i\varphi(x)\leftrightarrow x\in W_e\bigr),
\tag{S2}
\] with effective renaming of bound and free variables as needed. A
\emph{particular-code endorsement} by \(T\) is an actual theorem
\(T\vdash\operatorname{Code}_{i,\varphi}(\bar e)\) for a standard
numeral \(\bar e\). Merely proving an existential assertion that some
index exists is a different condition.

For the arithmetic results, \(\mathcal P_3\) is the recursive class of
designated prenex \(\Pi_3\) sentences, and \(\tau\) is its \(\Pi_3\)
partial truth predicate, false on other codes. The arithmetic theory \[
U_e=\mathsf I\Sigma_1+
\{\sigma\in\mathcal P_3:\ulcorner\sigma\urcorner\in W_e\}
\tag{S3}
\] uses enumeration witnesses in its proof predicate. The designated
\(\Pi_3\) sentence \(R_e\) is a fixed normalization of
\(\operatorname{RFN}_3(U_e)\). In particular, \[
\mathsf{PA}\vdash
\forall z(z\in W_e\rightarrow\tau(z))\rightarrow R_e.
\tag{S4}
\] The construction and proof of (S4) are {[}Main, Sections 4 and 5{]}.
This is a reflection statement for a fixed \(\Pi_3\) column over
\(\mathsf I\Sigma_1\), rather than reflection for an unrestricted
arithmetic theory.

The results are consequences of the same reduction, or standard
arithmetic arguments applied to it. Their inclusion does not assert that
they are new in the literature. In particular, the arithmetic
obstruction in Section 5 uses the usual second incompleteness theorem.

\section{2. Enumerating all later consequences from one
root}\label{enumerating-all-later-consequences-from-one-root}

\subsection{Proposition S1. Uniform enumeration of the full consequence
sets}\label{proposition-s1.-uniform-enumeration-of-the-full-consequence-sets}

Assume (S1), standard soundness, and particular-code endorsements \[
\forall j\geq1\ \forall\varphi(x)\ \exists e\quad
T_{j-1}\vdash\operatorname{Code}_{j,\varphi}(\bar e)
\tag{S5}
\] for all unary modal formulas. If the consequences of \(T_0\) are
recursively enumerable, then the relation \[
\{(j,\ulcorner\theta\urcorner):j\in\mathbb N,
\ \theta\text{ a sentence},\ T_j\vdash\theta\}
\tag{S6}
\] is recursively enumerable, uniformly from an enumeration of the root
consequences.

It is enough that the underlying consequence relation supports ordinary
logical conjunction and reflexivity of equality. The proposition does
not require arithmetic induction or known successor truthfulness.

\textbf{Proof.} The selection procedure \(B\) of {[}Main, Lemma 2{]}
uses only (S1), soundness, (S5) and the supplied root enumeration. Its
proof does not use \(\mathsf{PA}\) or successor truthfulness. It
returns, for every \(j\geq1\) and every unary formula \(\varphi\), an
index \(e\) such that \[
W_e=\{m:T_j\vdash\varphi(\bar m)\}.
\] For a sentence \(\theta\), take
\(\delta_\theta(x)=\theta\wedge x=x\). This has exactly the one free
variable \(x\), and ordinary logical deductions give \[
T_j\vdash\theta
\quad\Longleftrightarrow\quad
T_j\vdash\delta_\theta(0).
\tag{S7}
\] To semidecide the pair \((j,\ulcorner\theta\urcorner)\) for
\(j\geq1\), compute \(e=B(j,\ulcorner\delta_\theta\urcorner)\),
enumerate \(W_e\), and accept when 0 appears. The computation of \(B\)
terminates by its external induction. Acceptance is equivalent to the
right side of (S7). For \(j=0\), use the supplied root enumeration
directly. Dovetail these computations over all pairs to obtain an
enumeration of (S6). \hfill\(\square\)

This gives consequence sets, not the original axiom sets. It also does
not give the root a proof that the joint enumeration is total, or a
proof that it enumerates all the later machines correctly. Those
internal assertions were unnecessary in the main argument.

\section{3. Computable paths and finite
cycles}\label{computable-paths-and-finite-cycles}

\subsection{Proposition S2. No computable infinite endorsement
path}\label{proposition-s2.-no-computable-infinite-endorsement-path}

Let \(\{T_i:i\in\mathbb N\}\) be a family in the common recursive modal
language, satisfying (S1), standard soundness and ordinary arithmetic
\(\mathsf{PA}\). Let \(p:\mathbb N\to\mathbb N\) be total computable.
Suppose the consequences of \(T_{p(0)}\) are recursively enumerable.
There cannot be, at every \(n\), both \[
\forall\varphi(x)\ \exists e\quad
T_{p(n)}\vdash
\operatorname{Code}_{p(n+1),\varphi}(\bar e)
\tag{S8}
\] and \[
T_{p(n)}\vdash
\forall x(K_{p(n+1)}\tau(x)\rightarrow\tau(x)).
\tag{S9}
\] The function \(p\) need not be injective.

\textbf{Proof.} Adapt {[}Main, Lemma 2{]} with the natural number \(n\)
denoting position on the path. At the base, search the enumeration of
\(T_{p(0)}\) for a particular code of the \(K_{p(1)}\) column. At the
step, replace the literal operator name \(K_{j+1}\) in that lemma by
\(K_{p(n+1)}\). Since \(p\) is total computable, each required literal
name and each nested code formula can be computed. The same simultaneous
induction proves termination and retains the endorsement by the actual
preceding theory.

The wrapper programs of {[}Main, Section 5{]} now first compute the
required values of \(p\) and execute this adapted procedure. Their
indices are still obtainable uniformly by parameter compilation; no
computation of \(p\) has to be completed in order to compile a wrapper.
At each fixed position, its finite terminating computation is verified
in \(\mathsf I\Sigma_1\), just as in {[}Main, Lemma 6{]}.

Apply (S9) and (S4) to the endorsed \(\tau\)-column at each edge. The
proof in {[}Main, Section 7{]} gives a uniformly recursively enumerable
sequence of sound arithmetic extensions of \(\mathsf I\Sigma_1\), each
proving \(\Pi_3\) reflection for the next. This contradicts {[}Main,
Theorem 7{]}, the Pakhomov--Walsh non-descent theorem. Repetitions in
\(p\) affect neither the algorithm nor the contradiction.
\hfill\(\square\)

\subsection{Corollary S3. No finite
cycle}\label{corollary-s3.-no-finite-cycle}

In such a sound arithmetic family, make a directed edge \(i\to j\) mean
that \(T_i\) has all the particular-code endorsements for \(K_j\) and
proves its \(\tau\)-truthfulness. If every vertex's consequence set is
recursively enumerable, the directed graph has no finite cycle,
including a loop.

\textbf{Proof.} A finite cycle can be repeated periodically to give a
total computable function \(p\) satisfying (S8) and (S9). Proposition S2
excludes it. \hfill\(\square\)

A graph without computable infinite paths can still have a noncomputable
infinite path. Consequently neither statement proves that every such
endorsement relation is well-founded. This is the precise limitation
relevant to Alexander's ill-founded examples in \emph{Self-referential
theories}, Theorem 79; see {[}Main, Section 8.2{]}.

\section{4. A sentence separating two adjacent
theories}\label{a-sentence-separating-two-adjacent-theories}

\subsection{\texorpdfstring{Proposition S4. An endorsed edge gives
explicit \(\Pi_3\)
separation}{Proposition S4. An endorsed edge gives explicit \textbackslash Pi\_3 separation}}\label{proposition-s4.-an-endorsed-edge-gives-explicit-pi_3-separation}

Let \(S,T\) be sound theories containing ordinary \(\mathsf{PA}\), with
a canonical operator \(K_T\) for the consequences of \(T\). Suppose that
for a standard number \(e\), \[
S\vdash\forall x(K_T\tau(x)\leftrightarrow x\in W_{\bar e}),
\qquad
S\vdash\forall x(K_T\tau(x)\rightarrow\tau(x)).
\tag{S10}
\] Then the designated \(\Pi_3\) sentence \(R_e\) satisfies \[
S\vdash R_e,
\qquad
T\nvdash R_e.
\tag{S11}
\] In particular, the designated \(\Pi_3\) consequences of \(S\) are not
a subset of those of \(T\).

\textbf{Proof.} Soundness of \(S\), the first theorem of (S10), and
canonical semantics give \[
W_e=\{m:T\vdash\tau(\bar m)\}.
\tag{S12}
\] Thus {[}Main, Lemma 5{]} applies: \(U_e\) is sound and has exactly
the designated \(\Pi_3\) consequences of \(T\). Inside \(S\), the two
theorems of (S10) give \(\forall x(x\in W_e\rightarrow\tau(x))\).
Equation (S4) gives \(S\vdash R_e\).

Suppose \(T\vdash R_e\). Because \(R_e\) is a designated \(\Pi_3\)
sentence, equality of the \(\Pi_3\) consequence sets gives
\(U_e\vdash R_e\). The fixed base proves that uniform \(\Pi_3\)
reflection implies consistency: apply reflection to a fixed false
arithmetic sentence, whose falsehood the base proves. Consequently
\(U_e\) proves its own consistency for its specified witnessed proof
predicate. This contradicts the second incompleteness theorem, since
\(U_e\) is a consistent recursively enumerable extension of
\(\mathsf I\Sigma_1\) with the ordinary witnessed derivability
conditions. Therefore \(T\nvdash R_e\). \hfill\(\square\)

This gives a specific separating sentence once an endorsed code is
supplied. It does not prove the reverse inclusion of \(\Pi_3\)
consequence sets, assign an ordinal rank, or show that all arithmetic
consequences of one theory extend those of the other. The recursive
enumerability needed for incompleteness is that of \(U_e\); the original
modal theory \(T\) need not itself be supplied effectively.

\section{\texorpdfstring{5. The obstruction to true \(\Sigma_3\)
conditional
codes}{5. The obstruction to true \textbackslash Sigma\_3 conditional codes}}\label{the-obstruction-to-true-sigma_3-conditional-codes}

The final observation concerns a weaker-looking code assertion. Instead
of endorsing a particular code outright, a theory might prove that it is
correct under a true arithmetic condition. A low-complexity condition
does not evade the obstruction below when the two theories have the
specified inclusion of \(\Pi_3\) consequences.

\subsection{\texorpdfstring{Lemma S5. A true \(\Sigma_3\) guard cannot
imply own
reflection}{Lemma S5. A true \textbackslash Sigma\_3 guard cannot imply own reflection}}\label{lemma-s5.-a-true-sigma_3-guard-cannot-imply-own-reflection}

Let \(V\) be a sound recursively enumerable extension of
\(\mathsf I\Sigma_1\), with an ordinary witnessed proof predicate and
its designated uniform \(\Pi_3\) reflection sentence \(R_V\). If
\(\sigma\) is a true arithmetical \(\Sigma_3\) sentence, then \[
V\nvdash\sigma\rightarrow R_V.
\tag{S13}
\]

\textbf{Proof.} Suppose the indicated implication were a theorem. Put
\(V'=V+\sigma\). This is a sound, recursively enumerable extension of
\(\mathsf I\Sigma_1\), and it proves \(R_V\). Let \(\nu\) be a fixed
designated \(\Pi_3\) normalization of \(\neg\sigma\); the base proves
their equivalence.

The proof transformation for the deduction theorem, applied to the one
additional axiom \(\sigma\), is verified in \(\mathsf I\Sigma_1\). It
gives \[
\operatorname{Pr}_{V'}(\ulcorner\bot\urcorner)
\rightarrow
\operatorname{Pr}_V(\ulcorner\nu\urcorner).
\tag{S14}
\] In \(V'\), reflection \(R_V\), the truth biconditional for the fixed
sentence \(\nu\), and (S14) imply
\(\operatorname{Pr}_{V'}(\ulcorner\bot\urcorner)\rightarrow\nu\). But
\(V'\) also proves \(\sigma\), hence \(\neg\nu\). Thus \(V'\) proves
\(\operatorname{Con}(V')\), contradicting the second incompleteness
theorem. \hfill\(\square\)

Truth of \(\sigma\) ensures the consistency of \(V'\); the use of
\(\Sigma_3\) ensures that its negation is within the reflection class.
These roles should not be conflated. The statement does not extend to
arbitrary true \(\Pi_3\) guards: take \(\sigma=R_V\), which is true, and
the implication is a tautology.

\subsection{\texorpdfstring{Proposition S6. No true \(\Sigma_3\)-guarded
code under \(\Pi_3\)
inclusion}{Proposition S6. No true \textbackslash Sigma\_3-guarded code under \textbackslash Pi\_3 inclusion}}\label{proposition-s6.-no-true-sigma_3-guarded-code-under-pi_3-inclusion}

Let \(S,T\) be sound modal theories containing \(\mathsf{PA}\), with a
canonical operator \(K_T\). Suppose \[
\operatorname{Th}_{\mathcal P_3}(S)
\subseteq\operatorname{Th}_{\mathcal P_3}(T)
\tag{S15}
\] and \(S\vdash\forall x(K_T\tau(x)\rightarrow\tau(x))\). There are no
standard number \(e\) and true arithmetical \(\Sigma_3\) sentence
\(\sigma\) such that \[
S\vdash\sigma\rightarrow
\forall x(K_T\tau(x)\leftrightarrow x\in W_{\bar e}).
\tag{S16}
\]

\textbf{Proof.} If (S16) holds, its truth in the common standard
interpretation, together with the truth of \(\sigma\), gives the exact
equality (S12). Therefore \(U_e\) is sound and has precisely the
designated \(\Pi_3\) consequences of \(T\).

In \(S\), combine (S16) with its unconditional successor-truthfulness
theorem. Equation (S4) yields \[
S\vdash\sigma\rightarrow R_e.
\tag{S17}
\] Since \(\sigma\) is \(\Sigma_3\) and \(R_e\) is \(\Pi_3\), the
implication has a \(\Pi_3\) normalization provably equivalent in the
fixed base: write it as \(\neg\sigma\vee R_e\), and use closure of
\(\Pi_3\) under finite disjunction. Let \(\rho\in\mathcal P_3\) be that
normalization. By (S15), \(T\vdash\rho\). By equality of the designated
\(\Pi_3\) consequence sets, \(U_e\vdash\rho\), hence
\(U_e\vdash\sigma\rightarrow R_e\). This contradicts Lemma S5 for
\(V=U_e\). \hfill\(\square\)

This is a single-edge obstruction with an explicit inclusion hypothesis.
It assumes unconditional knowledge of the \(\tau\)-truthfulness clause;
it does not establish a corresponding result for conditional
truthfulness. It also does not rule out every conditional code, every
higher-complexity arithmetic guard, or an externally nonempty set of
code candidates without an endorsed numeral.

\section{6. Dependencies and status}\label{dependencies-and-status}

Proposition S1 uses only the effective selection argument of {[}Main,
Lemma 2{]}. Proposition S2 uses the full arithmetic reduction and its
published Pakhomov--Walsh input. Corollary S3 is a finite combinatorial
consequence of S2. Proposition S4 uses the column conversion and
ordinary second incompleteness. Lemma S5 is a direct finite-extension
argument using second incompleteness. Proposition S6 combines that lemma
with the column conversion and the explicitly assumed inclusion (S15).

These are hand proofs, with their assumptions displayed. They are not
part of the separate Lean formalization of the dissertation's Questions
3-1-1 through 3-1-3. No independent human review or priority claim is
asserted. The main paper contains the source locations, complete
arithmetic conventions and bibliography. For convenience, the two
principal external sources are Alexander's dissertation
\href{https://rave.ohiolink.edu/etdc/view?acc_num=osu1366362999}{OhioLINK
record} and Pakhomov--Walsh, \emph{Reflection ranks and ordinal
analysis},
\href{https://arxiv.org/abs/1805.02095v2}{arXiv:1805.02095v2}, Theorem
3.4.

\end{document}
